\let\iint\relax
\let\iiint\relax
%%% Работа с русским языком
\usepackage{cmap}					% поиск в PDF
\usepackage{mathtext} 				% русские буквы в формулах
\usepackage[T2A]{fontenc}			% кодировка
\usepackage[utf8]{inputenc}			% кодировка исходного текста
\usepackage[english,russian]{babel}	% локализация и переносы
\usepackage{indentfirst}            % красная строка в первом абзаце
\frenchspacing                      % равные пробелы между словами и предложениями
\usepackage{epigraph}
%%% Дополнительная работа с математикой
\usepackage{amsmath}
\usepackage{amsfonts,amssymb,amsthm,mathtools} % пакеты AMS
%\DeclareMathOperator{\arcctg}{arcctg}
\usepackage{icomma}                                    % "Умная" запятая
\usepackage{nccfloats}

%%% Свои символы и команды
\usepackage{centernot} % центрированное зачеркивание символа
\usepackage{stmaryrd}  % некоторые спецсимволы
%\usepackage{bm}

\usepackage{titletoc}
\usepackage{blindtext}

% Обычный формат секций
\titlecontents{section}
[0em]
{\vspace{1ex}}
{\thecontentslabel\quad}
{}
{\titlerule*[1pc]{.}\contentspage}

% Обычный формат подсекций
\titlecontents{subsection}
[2em]
{}
{\thecontentslabel\quad}
{}
{\titlerule*[1pc]{.}\contentspage}

% Команда для "продвинутых" секций
\makeatletter
\newcommand{\advsection}[1]{%
	\renewcommand{\thesection}{\arabic{section}*}% Временно добавляем звёздочку
	\section{#1}%
	\renewcommand{\thesection}{\arabic{section}}% Возвращаем стандартный формат
}
\makeatother

\makeatletter
\newcommand{\veryadvsection}[1]{%
	\renewcommand{\thesection}{\arabic{section}**}% Временно добавляем звёздочку
	\section{#1}%
	\renewcommand{\thesection}{\arabic{section}}% Возвращаем стандартный формат
}
\makeatother

% Команда для "продвинутых" подсекций
\makeatletter
\newcommand{\advsubsection}[1]{%
	\renewcommand{\thesubsection}{\arabic{section}.\arabic{subsection}* }% Временно добавляем звёздочку
	\subsection{#1}%
	\renewcommand{\thesubsection}{\arabic{section}.\arabic{subsection}}% Возвращаем стандартный формат
}
\makeatother

\makeatletter
\newcommand{\veryadvsubsection}[1]{%
	\renewcommand{\thesubsection}{\arabic{section}.\arabic{subsection}**}% Временно добавляем звёздочку
	\subsection{#1}%
	\renewcommand{\thesubsection}{\arabic{section}.\arabic{subsection}}% Возвращаем стандартный формат
}
\makeatother

\usepackage{comment} % позволяет использовать \begin{comment} глыба текста \end{comment}, чтобы закомментировать глыбу текста.

\usepackage{cancel} % cancel culture
%\usepackage{quiver}
\renewcommand{\epsilon}{\ensuremath{\varepsilon}}
\renewcommand{\phi}{\ensuremath{\varphi}}
\renewcommand{\kappa}{\ensuremath{\varkappa}}
\renewcommand{\le}{\ensuremath{\leqslant}}
\renewcommand{\leq}{\ensuremath{\leqslant}}
\renewcommand{\ge}{\ensuremath{\geqslant}}
\renewcommand{\geq}{\ensuremath{\geqslant}}
\renewcommand{\emptyset}{\ensuremath{\varnothing}}

\DeclareMathOperator{\sgn}{sgn}
\DeclareMathOperator{\Ker}{Ker}
\DeclareMathOperator{\cov}{cov}
\DeclareMathOperator{\Int}{Int}
\DeclareMathOperator{\im}{Im}
\DeclareMathOperator{\re}{Re}
\DeclareMathOperator{\rank}{rank}
\DeclareMathOperator{\grad}{grad}
\DeclareMathOperator{\ob}{Ob}
\DeclareMathOperator*{\astm}{\ast}
\DeclareMathOperator*{\sqcupm}{\sqcup}
\DeclareMathOperator{\id}{id}
\DeclareMathOperator{\Hess}{Hess}
\DeclareMathOperator{\tr}{tr}


\newcommand{\N}{\mathbb{N}}
\newcommand{\Z}{\mathbb{Z}}
\newcommand{\Q}{\mathbb{Q}}
\newcommand{\R}{\mathbb{R}}
\newcommand{\HH}{\mathbb{H}}
\newcommand{\Tor}{\mathbb{T}}
\newcommand{\Sph}{\mathbb{S}}
\newcommand{\RP}{\mathbb{RP}}
\newcommand{\D}{\mathcal{D}}
\newcommand{\Cm}{\mathbb{C}}
\newcommand{\F}{\mathbb{F}}
\newcommand{\Id}{\mathrm{Id}}
\newcommand{\DIF}{\mathcal{DIF}}
\newcommand{\Rim}{\mathcal{R}}
\newcommand{\M}{\mathcal{M}}
\newcommand{\G}{\mathcal{G}}
\newcommand{\B}{\mathcal{B}}
\newcommand{\X}{\mathcal{X}}
\renewcommand{\P}{\mathbb{P}}
\newcommand{\Ex}{\mathbb{E}}
\renewcommand{\D}{\mathbb{D}}


% привееет





\newcommand{\imp}[2]{
	(#1\,\,$\ra$\,\,#2)\,\,
}
\newcommand{\System}[1]{
	\left\{\begin{aligned}#1\end{aligned}\right.
}
\newcommand{\Root}[2]{
	\left\{\!\sqrt[#1]{#2}\right\}
}
\newcommand{\bt}[1]{
	\boldsymbol{#1}
}
\renewcommand\labelitemi{$\triangleright$}

\let\bs\backslash
\let\Lra\Leftrightarrow
\let\lra\leftrightarrow
\let\Ra\Rightarrow
\let\ra\rightarrow
\let\La\Leftarrow
\let\la\leftarrow
\let\emb\hookrightarrow
\let\rr\rightrightarrows
\newcommand{\nrr}{\bcancel{\rightrightarrows}}

\newcommand{\xrr}{\underset{X}{\rightrightarrows}}

\newcommand{\xnrr}{\underset{X}{\bcancel{\rightrightarrows}}}
%%% Перенос знаков в формулах (по Львовскому)
\newcommand{\hm}[1]{#1\nobreak\discretionary{}{\hbox{$\mathsurround=0pt #1$}}{}}

%%% Работа с картинками
\usepackage{graphicx}    % Для вставки рисунков
\setlength\fboxsep{3pt}  % Отступ рамки \fbox{} от рисунка
\setlength\fboxrule{1pt} % Толщина линий рамки \fbox{}
\usepackage{wrapfig}     % Обтекание рисунков текстом

%%% Работа с таблицами
\usepackage{array,tabularx,tabulary,booktabs} % Дополнительная работа с таблицами
\usepackage{longtable}                        % Длинные таблицы
\usepackage{multirow}                         % Слияние строк в таблице

%%% Теоремы
\theoremstyle{plain}
\newtheorem{theorem}{Теорема}[section]
\newtheorem{theoremdefinition}{Теорема-определение}[section]
\newtheorem{lemma}{Лемма}[section]
\newtheorem{proposition}{Утверждение}[section]
\newtheorem*{exercise}{Упражнение}
\newtheorem*{problem}{Задача}
\newtheorem*{question}{Вопрос}
\newtheorem*{answer}{Ответ}
\newtheorem*{explanation}{Пояснение}
\newtheorem*{fact}{Факт}

\theoremstyle{definition}
\newtheorem{definition}{Определение}[section]
\newtheorem*{corollary}{Следствие}
\newtheorem*{note}{Примечание}
\newtheorem*{reminder}{Напоминание}
\newtheorem*{example}{Пример}
\newtheorem*{examples}{Примеры}
\newtheorem*{counterexample}{Контрпример}
% \newtheorem*{solution}{Решение}
\newtheorem*{properties}{Свойства}
\newtheorem*{remark}{Замечание}

\theoremstyle{remark}
\newtheorem*{solution}{Решение}

%%% Оформление страницы
\usepackage{vmargin}
\usepackage{extsizes}     % Возможность сделать 14-й шрифт
%\usepackage{geometry}     % Простой способ задавать поля
\usepackage{setspace}     % Интерлиньяж
\usepackage{enumitem}     % Настройка окружений itemize и enumerate
\setlist{leftmargin=25pt} % Отступы в itemize и enumerate
%\geometry{top=25mm}    % Поля сверху страницы
%\geometry{bottom=30mm} % Поля снизу страницы
%\geometry{left=20mm}   % Поля слева страницы
%\geometry{right=20mm}  % Поля справа страницы
\setmarginsrb{2 cm}{11 mm}{2 cm}{14 mm}{0pt}{9mm}{0pt}{7mm}

\setlength\parindent{15pt}        % Устанавливает длину красной строки 15pt
\linespread{1.3}                  % Коэффициент межстрочного интервала
%\setlength{\parskip}{0.5em}      % Вертикальный интервал между абзацами
%\setcounter{secnumdepth}{0}      % Отключение нумерации разделов
%\setcounter{section}{-1}         % Нумерация секций с нуля
\usepackage{multicol}			  % Для текста в нескольких колонках
%\usepackage{soulutf8}             % Модификаторы начертания
\usepackage{soul}

%%% Содержаниие
\usepackage{tocloft}
\tocloftpagestyle{main}
%\setlength{\cftsecnumwidth}{2.3em}
%\renewcommand{\cftsecdotsep}{1}
%\renewcommand{\cftsecpresnum}{\hfill}
%\renewcommand{\cftsecaftersnum}{\quad}

\usepackage{biblatex}
\usepackage{csquotes}
%\addbibresource{sources.bib}

%%% Шаблонная информация для титульного листа
\newcommand{\CourseName}{Аналитическая механика}
\newcommand{\FullCourseNameFirstPart}{\so{Аналитическая механика}}
\newcommand{\SemesterNumber}{-}
\newcommand{\LecturerInitials}{Д. А. Притыкин}
\newcommand{\CourseDate}{\today}
\newcommand{\AuthorInitials}{Бадоля Пётр}
\newcommand{\VKLink}{https://vk.com/akzium}
\newcommand{\TGLink}{https://t.me/Postscripter}
\newcommand{\GithubLink}{https://github.com/daniild71r/lectures_tex_club}


%%% Колонтитулы
\usepackage{titleps}
\newpagestyle{main}{
	\setheadrule{0.4pt}
	\sethead{\CourseName}{}{\hyperlink{intro}{\;Назад к содержанию}}
	\setfootrule{0.4pt}                       
	\setfoot{ФПМИ МФТИ, \CourseDate}{}{\thepage} 
}
\pagestyle{main}  

%%% Нумерация уравнений
\makeatletter
\def\eqref{\@ifstar\@eqref\@@eqref}
\def\@eqref#1{\textup{\tagform@{\ref*{#1}}}}
\def\@@eqref#1{\textup{\tagform@{\ref{#1}}}}
\makeatother                      % \eqref* без гиперссылки
\numberwithin{equation}{section}  % Нумерация вида (номер_секции).(номер_уравнения)
\mathtoolsset{showonlyrefs=false} % Номера только у формул с \eqref{} в тексте.

%%% Гиперссылки
\usepackage[unicode, pdftex]{hyperref}
%\usepackage[usenames,dvipsnames,svgnames,table,rgb]{xcolor}
\usepackage[dvipsnames,svgnames,table,rgb]{xcolor}
\hypersetup{
	unicode=true,            % русские буквы в раздела PDF
	colorlinks=true,       	 % Цветные ссылки вместо ссылок в рамках
	linkcolor=black!15!blue, % Внутренние ссылки
	citecolor=green,         % Ссылки на библиографию
	filecolor=magenta,       % Ссылки на файлы
	urlcolor=NavyBlue,       % Ссылки на URL
}

%%% Графика
\usepackage{tikz}        % Графический пакет tikz
\usepackage{pgfplots}
\pgfplotsset{width=10cm,compat=1.9} % Графики
\usepackage{tikz-cd}     % Коммутативные диаграммы
\usepackage{tkz-euclide} % Геометрия
\usepackage{stackengine} % Многострочные тексты в картинках
\usetikzlibrary{patterns}
\usetikzlibrary{angles, babel, quotes}

%\includeonly{lectures/lect05,lectures/lect06}  % Скомпилировать только часть лекций


\everymath{\displaystyle} % Чтобы формулы большие были


\usepackage{mathrsfs}
\usepackage{tikz-3dplot}
\tdplotsetmaincoords{70}{110}